\section{Limitations}
\label{sec:limitations}

\paragraph{Held-out size.} 429 test frames over 97 scenes. The oracle and
ceiling analyses of \cref{sec:bound} use 27 and 51 frames respectively, one per
scene; that sample runs approximately \SI{1.15}{\decibel} pessimistic on clean
against the full set (\SI{-4.15}{\decibel} at $\alpha = 1$ where the full 429
read \SI{-3.00}{\decibel}). Ordering is reliable; absolute values will move.

\paragraph{Distribution shift.} The train/test shift of \cref{sec:shift} is a
confound in the clean result. \Cref{sec:ceiling} argues it is not the binding
constraint, but we have not retrained under a matched sampler to prove it.

\paragraph{The oracle is an upper bound.} It reads the ground truth.

\paragraph{Single-image only.} The temporal refiner is not evaluated here; its
held-out set is 4 validation and 5 test clips of one scene each, too small to
report.

\paragraph{The analytic baseline is privileged on clean corpus frames.}
\Cref{eq:baseline} is the exact inverse of the transform that rendered the
corpus SDR side, so on clean frames it is being handed the forward pipeline. It
is an unusually strong control, which is what makes the comparison sharp, and it
is also unusually strong in a way no baseline would be on footage from elsewhere.
A clean-condition regression against it should be read as ``worse than knowing
the exact forward transform'', not as a general statement about learned inverse
tone mapping against whatever baseline a real pipeline would have. The degraded
condition does not have this property, since no analytic inverse knows the
degradation.

\paragraph{One degradation model, and a gate supervised on it.} \textsc{hard}
is our own seeded camera and codec pipeline, not a corpus of real degraded
footage. This bears harder on the gate than on the model: the gate is trained by
binary cross-entropy against \emph{our} degradation label, so what it learns may
be the signature of this generator, JPEG blocking, 4:2:0 chroma, the particular
banding and tone-curve family, rather than the property we want, which is that a
frame's shadows need reconstruction. Nothing in our evaluation separates the
two, because the test condition is generated the same way as the training
condition. The experiment that would separate them is to hold out a degradation
family entirely, train on the rest and test on it, and better still to score
real degraded footage that our pipeline never touched. Until that is run, the
gate's 75.5\% should be read as accuracy on this degradation model.

\paragraph{One published method, not four.} ExpandNet is scored on our split
(\cref{sec:expandnet}); Santos et al., Eilertsen et al.~and Liu et al.~have
public code and are not. One comparison answers ``compared to what?'' and does
not rank the field, and the caveats of \cref{sec:expandnet}, namely
training-domain mismatch, an unclamped reference and a per-frame exposure fit,
apply to the one row we have.

\paragraph{No claim here rests on a reimplementation of prior work.}
\Cref{sec:related} places the result among families of methods and cites them
bibliographically; it does not reproduce their numbers, and we did not run their
code.

\paragraph{Three seeds is a spread, not a distribution.} The gate is positive on
all four measures in all three runs, but $n = 3$ supports ``the sign is stable'',
not a confidence interval. The clean PU21 spread ($+0.07$ to $+0.71$) is wide
relative to its mean.

\paragraph{Deployment carries the gate as a scalar only.} The interactive viewer
receives the predicted shadow weight per frame and multiplies the shadow prior
by it in the compositing shader; CPU/GPU parity holds to
$8.4 \times 10^{-6}$ relative error across 11 cases spanning $\sw = 0$ to
$\sw = 1$. The weight cannot be folded into the residual the way a scalar
residual scale can, because it enters inside the $\max$ of \cref{eq:prior} and
is not linear in the residual.
